\section{Tasks are too big}
\label{section:node-level:task-size-too-big}

% If tasks (or parallel sections) are too small, we pay an excessive overhead.
% The concurrency level might be brilliant, but that does not help us if we have too much task creation and task clean-up overhead.
% 
% \begin{hypothesis}
%  There are enough tasks and they are not excessively synchronising each other, but they are just too small and hence create too much administration overhead.
% \end{hypothesis}
% 
% 
% \noindent
% Too many tiny tasks will manifest in a lot of time spent in the actual tasking runtime.
% Differen to excessive synchronisation (Section~\ref{section:node-level:synchronisation}), we however barely will see locks or spinning unless we start to profile our task runtime itself.
% 
% \paragraph{Data collection}
% 
% \begin{instructions}{\otter}
%  to be written
% \end{instructions}
% 
% 
% \begin{instructions}{\vtune}
%  to be written
% \end{instructions}
% 
% 
% \paragraph{Evaluation}
% 
% The extreme case of tiny tasks are empty tasks which some codes use to model dependencies, i.e.~to synchronise sets of tasks in a task graph.
% Another reason for tiny tasks is the classic CUDA-style mapping, where a task loop is bigger than the number of tasks and the loop body is then embedded into one big \texttt{if} statement.
% 
% 
% 
% \paragraph{Optimisation and follow-up steps}
% 
% There are two different strategies to get rid of too small tasks.
% On the one hand, we can try to avoid that small tasks are created in the first place.
% For data-parallel regions (e.g.~OpenMP's \texttt{parallel for}), we can use the minimal chunk/grain size to ensure that a loop is not cut down into tiny segments.
% On the other hand, we can stick to small tasks and fuse them later into bigger ones.
% This approach is also sometimes called batching (in the linear algebra world).
% 
% 
% This second approach usually requires some more coding work, i.e.~while it is easy to restrict the chunk size of a parallel loop, it is not trivial to fuse tasks once they are generated. 
% Actually, no mainstream tasking backend supports this.
% Our code has to provide means to package multiple logical tasks into one physical task.
% 
% 
% Once we alter the grain sizes, it is mandatory to return to the overall assessment.
% It is not unlikely that we harm the balancing (Section~\ref{section:node-level:balancing}) and have to balance carefully between a good balancing and concurrency level on the side and limited overhead o the other.
% 
% 
% \paragraph*{How it works}
% 
% \begin{itemize}[noitemsep]
%   \item[$\boxplus$] Validate that the administrative (non-science) fraction of the code scales with the scientific core calculations.
% \end{itemize}
% 
% 
% \paragraph*{How it does not work}
% 
% \begin{itemize}[noitemsep]
%   \item[$\boxminus$] The administrative overhead will always increase relatively, so an increase in itself is not a flaw. An extraordinary increase coupled with limited scalability however is.
% \end{itemize}
